\documentclass[11pt]{article}
\usepackage[margin=1in]{geometry}
\usepackage{amsmath,amssymb,amsthm,booktabs,array,microtype}
\usepackage[hidelinks]{hyperref}
\newtheorem{theorem}{Theorem}[section]
\newtheorem{lemma}[theorem]{Lemma}
\newtheorem{corollary}[theorem]{Corollary}
\newcommand{\RP}{\mathbb{RP}}
\title{Projective-plane triangulations yield\\orderable cographic matroids}
\author{Carptopus\\\small\texttt{carptopus@163.com}}
\date{Version 0.1-beta, 30 August 2026}

\begin{document}
\maketitle

\begin{abstract}
Crenshaw and Oxley conjectured that every 3-connected orderable binary matroid
is graphic. We give a topological construction that disproves this conjecture.
If $T$ is a finite simplicial triangulation of the real projective plane and
$G$ is its 1-skeleton, then the cographic matroid $M^*(G)$ is orderable. The
ordering is induced by face corners. For every bond of $G$, the resulting
2-regular adjacency graph is identified with a piecewise-linear separating
level set. An Euler-characteristic argument shows that this level set has one
component. Starting from the six-vertex triangulation with 1-skeleton $K_6$
and repeatedly applying stellar subdivision to triangular faces gives
infinitely many 3-connected, regular, non-graphic, orderable binary matroids.
We also give the explicit $K_6$ certificate and an exact finite checker as a
reproducibility aid.
\end{abstract}

\section{Introduction}

An ordering of a matroid assigns a reversible cyclic ordering to every
circuit. It is \emph{consistent} when common adjacent pairs are treated
compatibly across the circuit family; a matroid admitting such an assignment
is \emph{orderable}~\cite{crenshaw-oxley}. Crenshaw and Oxley proved strong
structural restrictions on orderable matroids and conjectured that every
3-connected orderable binary matroid is graphic
\cite[Conjecture~4]{crenshaw-oxley}. Their theorem for 4-connected regular
matroids leaves the exact 3-connected boundary open.

We show that this boundary contains a natural infinite family of
counterexamples. The triangular faces of a projective-plane triangulation
define a single global adjacency relation on the edges of its 1-skeleton. On
each bond this relation has degree two. The main point is connectedness. The
bond can be realized as the mid-level set of a two-valued piecewise-linear
function on the surface. Because both sides of a bond are connected, the two
complementary subsurfaces are connected. The positive Euler characteristic of
the real projective plane then forces the common boundary to be one circle.

The six-vertex triangulation makes the smallest member especially concrete:
its 1-skeleton is $K_6$, so $M^*(K_6)$ itself is an orderable, regular, binary,
3-connected, non-graphic matroid. Repeated face subdivisions produce the
infinite family.

\section{Definitions and a global-adjacency criterion}

All graphs and simplicial complexes are finite. For a connected graph $G$, the
circuits of the cographic matroid $M^*(G)$ are the bonds of $G$, that is, the
inclusion-minimal nonempty edge cuts. For a nonempty proper vertex set
$S\subset V(G)$, the cut $\delta(S)$ is a bond exactly when both induced
subgraphs $G[S]$ and $G[V(G)\setminus S]$ are connected.

We use the following immediate sufficient form of consistency.

\begin{lemma}[Global adjacency]\label{lem:global}
Let $M$ be a matroid. Suppose that a simple graph $\Lambda$ on $E(M)$ has the
property that, for every circuit $C$ of $M$, the induced graph $\Lambda[C]$ is
a single cycle. Reading that cycle in either direction gives consistent
reversible cyclic orderings of all circuits. In particular, $M$ is orderable.
\end{lemma}

\begin{proof}
All circuit orderings are restrictions of the same unordered adjacency
relation. Hence an adjacent pair shared by two circuits is adjacent in both
induced cycles, which is the required consistency.
\end{proof}

\section{The projective-plane theorem}

\begin{theorem}\label{thm:rp2}
Let $T$ be a finite simplicial triangulation of $\RP^2$, and let
$G=T^{(1)}$ be its 1-skeleton. Then $M^*(G)$ is orderable.
\end{theorem}

\begin{proof}
Define a simple graph $\Lambda_T$ on $E(G)$ as follows. Two edges of $G$ are
adjacent in $\Lambda_T$ when they are distinct edges of a common triangular
face of $T$. Since $T$ is simplicial, a pair of incident edges belongs to at
most one common face, so this is a well-defined global adjacency relation.

Fix a bond $B=\delta(S)$. Every edge $uv\in B$ belongs to exactly two
triangular faces. In a face $uvw$, the third vertex $w$ lies on exactly one
side of the cut. Consequently exactly one of $uw,vw$ is another edge of $B$,
and that edge is adjacent to $uv$ in $\Lambda_T$. The two faces incident with
$uv$ have different third vertices, so they supply two different neighbors.
Thus $\Lambda_T[B]$ is 2-regular and hence is a disjoint union of cycles.

It remains to show that there is only one cycle. Define a piecewise-linear map
\[
h:|T|\longrightarrow[0,1]
\]
by setting $h(v)=0$ for $v\in S$, $h(v)=1$ for $v\notin S$, and extending
affinely over every simplex. Put
\[
C=h^{-1}(1/2),\qquad
N_0=h^{-1}([0,1/2]),\qquad
N_1=h^{-1}([1/2,1]).
\]
In every mixed triangular face, $C$ is the segment joining the midpoints of
its two crossing edges. These segments join across common edges and form a
disjoint union of simple closed curves. Their combinatorial adjacency graph is
precisely $\Lambda_T[B]$. Since $1/2$ is not the value of any vertex, the same
triangle-by-triangle local model shows that $N_0,N_1$ are compact surfaces and
that $C=\partial N_0=\partial N_1$.

The subspace $N_0$ strongly deformation retracts onto the full induced
subcomplex $T[S]$. Indeed, in each simplex let $p$ be the sum of the
barycentric coordinates at vertices in $S$. On $N_0$ we have $p\ge 1/2$.
Join each point by the straight-line homotopy to the point obtained by dividing
its $S$-coordinates by $p$ and setting the remaining coordinates to zero.
Along this homotopy the sum of the $S$-coordinates is
$(1-t)p+t\ge1/2$, so the homotopy stays in $N_0$; the simplexwise formulas
agree on common faces. The same argument applies to $N_1$ and
$T[V(G)\setminus S]$. Since $B$ is a bond, the
1-skeleta $G[S]$ and $G[V(G)\setminus S]$ are connected. Therefore $N_0$ and
$N_1$ are connected compact surfaces with common boundary $C$.

Let $b$ be the number of components of $C$. A connected compact surface with
$b$ boundary components has Euler characteristic at most $2-b$, in both the
orientable and nonorientable cases. If $b\ge2$, then $\chi(N_0)\le0$ and
$\chi(N_1)\le0$. Because $C$ is a union of circles, $\chi(C)=0$, and additivity
gives
\[
1=\chi(\RP^2)=\chi(N_0)+\chi(N_1)-\chi(C)\le0,
\]
a contradiction. Hence $b=1$. Thus $\Lambda_T[B]$ is one cycle for every bond
$B$, and Lemma~\ref{lem:global} proves that $M^*(G)$ is orderable.
\end{proof}

\section{An infinite family of counterexamples}

On vertex set $\{0,1,2,3,4,5\}$, take the ten faces
\[
013,014,024,025,035,123,125,145,234,345.
\]
Every edge occurs in exactly two faces, every vertex link is a 5-cycle, and
$6-15+10=1$. This is a triangulation of $\RP^2$ whose 1-skeleton is $K_6$.

\begin{corollary}\label{cor:k6}
The matroid $M^*(K_6)$ is 3-connected, regular, binary, non-graphic, and
orderable. Hence it is a counterexample to Crenshaw--Oxley Conjecture~4.
\end{corollary}

\begin{proof}
Theorem~\ref{thm:rp2} gives orderability. Cographic matroids are regular and
binary. The simple 3-connected graph $K_6$ has a 3-connected cycle matroid,
and duality preserves matroid 3-connectivity. Finally, a cographic matroid of
a 3-connected graph is graphic only when the graph is planar; $K_6$ is
nonplanar~\cite{oxley}.
\end{proof}

Start with this triangulation $T_0$. Given a triangular face $abc$, insert a
new vertex $x$ in its interior and replace $abc$ by $abx,bcx,cax$. Denote the
resulting stellar subdivision by $T_{m+1}$, and write $G_m=T_m^{(1)}$.

\begin{theorem}\label{thm:family}
The matroids $M^*(G_m)$, $m\ge0$, are infinitely many pairwise nonisomorphic
3-connected, regular, binary, non-graphic, orderable matroids. Every member is
a counterexample to Crenshaw--Oxley Conjecture~4.
\end{theorem}

\begin{proof}
Stellar subdivision preserves the underlying surface and the simplicial
triangulation, so Theorem~\ref{thm:rp2} applies at every stage. We prove
3-connectivity of $G_m$ by induction. The new vertex $x$ has the three
neighbors $a,b,c$, while every old edge is retained. After deleting at most
two vertices, the surviving old graph is connected by the induction
hypothesis. If $x$ survives, at least one of $a,b,c$ also survives, so $x$ lies
in the same component. Thus $G_{m+1}$ is 3-connected.

Every $G_m$ contains the initial $K_6$ as a subgraph, and is therefore
nonplanar. The standard graphic--cographic planarity criterion, together with
3-connectivity, implies that $M^*(G_m)$ is non-graphic. It is cographic and
hence regular and binary, and its matroid 3-connectivity follows from that of
$G_m$. Finally, each subdivision adds three edges, so the ground-set sizes are
distinct.
\end{proof}

\section{The explicit $K_6$ certificate and verification boundary}

For reproducibility, the six vertex links, written as cyclic orders of the
other endpoints, are
\[
\begin{array}{c|c}
0&(1,3,5,2,4)\\
1&(0,3,2,5,4)\\
2&(0,4,3,1,5)\\
3&(0,1,2,4,5)\\
4&(0,1,5,3,2)\\
5&(0,2,1,4,3).
\end{array}
\]
Declaring two $K_6$-edges adjacent when they occur at a common face corner
gives 30 global adjacent pairs. The 31 bonds of $K_6$ split by cut type as six
of size 5, fifteen of size 8, and ten of size 9. The accompanying checker
reconstructs the binary cographic representation, all circuits, the global
adjacency relation, and matroid connectivity, and verifies that every induced
bond adjacency is a single cycle. A second checker performs stellar
subdivisions through a finite calibration range.

These computations are not used to prove the general theorem or the infinite
family. They are exact positive controls for the smallest example and the
implementation. No timeout or failed search is used as mathematical evidence.

\section{Scope and prior-work boundary}

The result disproves the stated 3-connected binary conjecture, but it does not
conflict with the 4-connected regular theorem of Crenshaw and Oxley. Every
$G_m$ retains triangles of the initial $K_6$; such a triangle is a
three-element cocircuit of $M^*(G_m)$, and together with 3-connectivity this
gives an exact 3-separation. We do not classify all
orderable cographic matroids or all surface triangulations that yield them.
The Euler-characteristic argument is specific to the positive Euler
characteristic of $\RP^2$; it does not assert an analogous theorem on the
torus or on higher-genus surfaces.

The literature search checked the published and arXiv versions of
\cite{crenshaw-oxley} and searched directly for the $M^*(K_6)$, cographic,
orderable, and projective-plane-triangulation combination. No direct prior
statement of the theorem or counterexample was located. This is a documented
search boundary, not a claim of exhaustive worldwide priority. The manuscript
is an internally verified candidate proof and has not undergone external peer
review.

\section*{AI-assisted research disclosure}

Carptopus is the responsible author and takes responsibility for the claims,
proofs, source use, and released artifacts. OpenAI Codex was used for
AI-assisted literature organization, proof development and stress testing,
exact verification, adversarial auditing, and manuscript preparation.

\begin{thebibliography}{9}
\bibitem{crenshaw-oxley}
S. Crenshaw and J. Oxley,
\emph{Ordering circuits of matroids},
Electronic Journal of Combinatorics 29 (2022), no.~4, Paper P4.31.
\url{https://doi.org/10.37236/11117}.

\bibitem{oxley}
J. Oxley,
\emph{Matroid Theory}, second edition,
Oxford University Press, 2011.

\bibitem{mohar-thomassen}
B. Mohar and C. Thomassen,
\emph{Graphs on Surfaces},
Johns Hopkins University Press, 2001.
\end{thebibliography}

\end{document}
